% Tahap prapemrosesan (kotak Preprocessing pada Gambar IV.1), digambar ulang dengan label yang sama.
% Gaya bentuk dan warna ada di diagram_style.tex.
\begin{tikzpicture}[dgbase]
\hyphenpenalty=10000 \exhyphenpenalty=10000 % jangan memenggal kata di dalam node
\node[dgdata] (ds)  at (0,0)     {Dataset};
\node[dgproc, minimum width=3.0cm, text width=2.7cm] (ubah) at (7.4,0) {Ubah label teks pada train menjadi vektor};
\node[dgdata, minimum width=2.8cm, text width=2.0cm] (bersih) at (11.3,0) {Dataset bersih (EEG $\to$ Vektor \textit{embedding})};
\node[dgproc, minimum width=2.6cm, text width=2.3cm] (std) at (3.2,-2.6) {Standardize Channel};
\node[dgproc, minimum width=2.6cm, text width=2.3cm] (lpf) at (7.4,-2.6) {Low Pass Filter};

\draw[dgarr] (ds) -- (ubah);
\draw[dgarr] (ubah) -- (bersih);
\draw[dgdash] (ds.south) |- (std.west);
\draw[dgdash] (std) -- (lpf);
\draw[dgdash] (lpf.north) -- (ubah.south);
\node[dglabel, align=left, anchor=west] at ($(ds.south)+(0.12,-0.55)$) {Pakai dataset\\lebih dari satu};
\end{tikzpicture}
